\documentclass[jisng5i.tex]{subfiles}
%subfile: ch_2.tex
\begin{document}

\section{Shrinkable and Normal Spaces} \label{sec:5-2}

Recall Definition \ref{def:1.1.2} in
{\definition{\label{def:5.2.1}}
Let $\mathscr{U} = \{U_\alpha: \alpha \in \Delta\}$ and $\mathscr{B} = \{B_\alpha: \alpha \in \Delta\}$ 
be covers of $X$. $\mathscr{B}$ is a 
\uwave{shrink}\index{shrink| \S\ref{def:5.2.1}\hfill} 
of $\mathscr{U}$ if $\overline{B_\alpha} \subset U_\alpha$ for each $\alpha\in \Delta$.

$X$ is 
\uwave{$\lambda$-shrinkable}\index{$\lambda$-shrinkable| \S\ref{def:5.2.1}\hfill} 
\index{shrinkable| \S\ref{def:5.2.1}\hfill} 
 every open $\lambda$-cover of $X$ has an open shrink.
}

{\lemma{\label{lem:5.2.2}}
\textup{(i)} A countably shrinkable space is normal and countable paracompact.

\textup{(ii)} For any $X$ the following are equivalent.

$\quad$ \textup{(a)} $X$ is $\lambda$-shrinkable.

$\quad$ \textup{(b)} Every open $\lambda$-cover of $X$ has a closed shrink.

$\quad$ \textup{(c)} Every open $\lambda$-cover of $X$ has a functionally open shrink.
}

\begin{proof}
(i) Suppose $X$ is countably shrinkable. Then $X$ is normal. By Theorem \ref{th:3.3.2}, 
$X$ is countably paracompact.

(ii) (a) $\Rightarrow$ (b) and (c) $\Rightarrow $(a). Clear.

(b) $\Rightarrow$ (c).
Assume (b). It is easy to show that $X$ is normal. Let 
$\mathscr{U} = \{U_\alpha: \alpha \in \Lambda\}$ be an open cover of $X$, and 
let $\{F_\alpha: \alpha \in \Lambda\}$ be a closed cover of $X$ 
such that $F_\alpha \subset U_\alpha$ for each $\alpha \in \Lambda$. 
Let $H_\alpha$ be a functionally open set such that 
$F_\alpha \subset H_\alpha \subset \overline{H_\alpha} \subset U_\alpha$. 
Then $\mathscr{H} = \{H_\alpha\}_{\alpha \in \Lambda}$ be a functionally open 
shrink of $\mathscr{U}$.
\end{proof}

{\definition{} \label{def:5.2.3}
\textup{(i)}  An open cover \(\mathscr{U} = \{ U_\alpha : \alpha \in \Delta \}\) of $X$ is a 
\uwave{$C\omega^2$-cover}\index{$C\omega^2$-cover| \S\ref{def:5.2.3}\hfill} 
if \(\mathscr{U}\) has a refinement \(\mathscr{F} = \bigcup_{m,n<\omega} \mathscr{F}_{m,n}\) 
where each \(\mathscr{F}_{m,n} = \{ F_{m,n,\alpha} : \alpha \in \Delta \}\) is a family of 
closed subsets of the subspace  \(X\setminus \cup \{\bigcup \mathscr{F}_{p, q}: (p, q) < (m,n) \}\)
and $F_{m, n, \alpha} \subset U_\alpha$ for all \(\alpha \in \Delta\).
The family \(\mathscr{F}\) is called a 
\uwave{$C\omega^2$-refinement}\index{$C\omega^2$-refinement| \S\ref{def:5.2.3}\hfill} 
of \(\mathscr{U}\).
\(X\) is 
\uwave{$\lambda$-$C\omega^2$-refinable}
\index{$\lambda$-$C\omega^2$-refinable| \S\ref{def:5.2.3}\hfill} 
\index{$C\omega^2$-refinable| \S\ref{def:5.2.3}\hfill} 
if every open \(\lambda\)-cover of \(X\) is a $C\omega^2$-cover.

\textup{(ii)}  An open cover \(\mathscr{U}\) is a 
\uwave{$C^*\omega^2$-cover}
\index{$C^*\omega^2$-cover| \S\ref{def:5.2.3}\hfill} 
 if \(\mathscr{U}\) has a $C\omega^2$-refinement 
 \(\mathscr{F} = \bigcup_{m,n<\omega} \mathscr{F}_{m,n}\) 
 such that each \(\bigcup \mathscr{F}_{m,n} \) is a closed subset of 
 $X\setminus \bigcup\{\mathscr{F}_{p, q}: (p, q) < (m, n)\}$.
}

{\theorem{} \label{th:5.2.4}
\(X\) is $\lambda$-shrinkable iff \(X\) is normal and $\lambda$-$C\omega^2$-refinable.
}

\begin{proof}
\((\Rightarrow)\) Suppose \(X\) is $\lambda$-shrinkable, then \(X\) is normal. 
Let \(\mathscr{U} = \{ U_\alpha : \alpha \in \Delta \}\) be an open cover of \(X\) and 
\(\mathscr{F}\) be a closed shrink of \(\mathscr{U}\). Let $\mathscr{F}_{0,0} = \mathscr{F}$ and 
 for each $(m, n) > (0,0), \mathscr{F}_{m,n}=\{\emptyset\}$.
 Then \(\bigcup_{m,n<\omega} \mathscr{F}_{m,n}\) is a \(C\omega^2\)-refinement of 
 \(\mathscr{U}\), so \(X\) is \(C\omega^2\)-refinable.
	
\((\Leftarrow)\) Suppose \(X\) is normal and \(C\omega^2\)-refinable. 
We first establish the following:
	
\textsc{Claim 1}. For any open \(\lambda\)-cover \(\mathscr{U} = \{U_\alpha: \alpha < \lambda\}\) 
of \(X\), there exists a sequence \(\{\mathscr{V}_{m,n} \}_{m,n<\omega}\) of 
\(\lambda\)-families of open 
subsets of $X$, 
where each \(\mathscr{V}_{m,n} = \{ V_{m,n,\alpha} : \alpha < \lambda\}\) 
such that \(\overline{V}_{m,n,\alpha} \subset U_\alpha\) for all \(\alpha < \lambda\) and 
\(\mathscr{V} = \bigcup_{m,n<\omega} \mathscr{V}_{m,n}\) covers \(X\).

\textsc{Proof of Claim 1}. We use induction on the well-ordered set 
$(\omega \times \omega, <)$. 
Let \(\mathscr{U} = \{ U_\alpha : \alpha < \lambda \}\) be an open \(\lambda\)-cover of \(X\) 
and let \(\mathscr{F} = \bigcup_{m,n<\omega} \mathscr{F}_{m,n}\) be a 
\(C\omega^2\)-refinement of \(\mathscr{U}\), where each 
\(\mathscr{F}_{m,n} = \{ F_{m,n,\alpha} : \alpha < \lambda\}\) is a family 
of closed subsets of the subspace 
\(X \setminus \cup\{\bigcup\mathscr{F}_{p,q}: (p, q) <  m,n\}\) with 
\(F_{m,n,\alpha} \subset U_\alpha\) for all \(\alpha < \lambda\).
	Let \(\mathscr{F}_{0,0} = \{\emptyset\}\), 
\(\mathscr{V}_{0,0} = \{\emptyset\}\) and
 \(\langle m,n \rangle \in \omega \times \omega\). 
 Assume that the $\lambda$-families 
 \(\mathscr{V}_{p,q} = \{ V_{p,q,\alpha} : \alpha < \lambda \}\) 
 of open sets  have been constructed for all \((p,q) <  ( m,n)\) to satisfy the conditions:
\begin{gather}
\overline{V}_{p,q,\alpha} \subset U_\alpha \text{ for all } \alpha < \lambda, \label{eq:5.2.4.1}\\
\bigcup_{(i, j) \le (p,q)} (\bigcup \mathscr{F}_{i,j}) \subset\bigcup_{(i, j) \le (p,q)} (\bigcup \mathscr{V}_{i,j}). \label{eq:5.2.4.2}
\end{gather}
Then we have
\begin{gather*}
\bigcup_{(p,q) < (m,n)} (\bigcup \mathscr{F}_{p,q}) \subset\bigcup_{(p,q) < (m,n)} (\bigcup \mathscr{V}_{p, q}). 
\end{gather*}
Let \(L = \bigcup_{(p,q) < (m,n)} (\bigcup \mathscr{V}_{p,q})\). 
Since each \(F_{m,n,\alpha} \setminus L\) is a closed subset of \(X\), there exists an open set \(V_{m,n,\alpha}\) of \(X\) such that
$$F_{m,n,\alpha} \setminus L \subset V_{m,n,\alpha} \subset \overline{V}_{m,n,\alpha} \subset U_\alpha.$$
Let \(\mathscr{V}_{m,n} = \{ V_{m,n,\alpha} : \alpha < \lambda \}\).  Then
\begin{gather*}
\bigcup \mathscr{F}_{m,n}\setminus L \subset \bigcup \mathscr{V}_{m,n}, \\
\bigcup_{(p,q) \le (m,n)} (\bigcup \mathscr{F}_{p,q}) \subset 
L\bigcup(\bigcup \mathscr{F}_{m,n})\subset
L\bigcup(\bigcup \mathscr{F}_{m,n}\setminus L)\subset
\bigcup_{(p,q) \le (m,n)} (\bigcup \mathscr{V}_{p,q}) 
\end{gather*}
Hence the open families \(\mathscr{V}_{m,n}\) satisfy 
Conditions \eqref{eq:5.2.4.1} and \eqref{eq:5.2.4.2} . 
By \eqref{eq:5.2.4.2}, \(\mathscr{V}\) covers \(X\). \textsc{Claim 1} is true.
	
\textsc{Claim 2.} \(X\) is countably paracompact.

\textsc{Proof of Claim 2}. Suppose $\mathscr{U} = \{U_k : k \in \omega\}$ be a 
countable increasing open cover of $X$. By \textsc{Claim 1}, there exists a sequence 
$\{\mathscr{V}_{m,n} : m, n < \omega\}$ of $\omega$-families of open sets of $X$ 
where each
$$V_{m,n} = \{\mathscr{V}_{m,n,k}: k < \omega\}$$
such that $\overline{V}_{m,n,k} \subset U_k$ for all $k$, and 
$\mathscr{V} = \bigcup_{m,n < \omega} \mathscr{V}_{m,n}$ covers $X$.
For $m < \omega$, let $G_m = \bigcup_{i,j,k=0}^m V_{i,j,k}$.
Since $\mathscr{U}$ is increasing, $\overline{G}_m \subset U_m$ and  
$\bigcup_{m < \omega} G_m = X$. Then
$\{G_k: k < \omega\}$ is an open shrink of $\mathscr{U}$. By Theorem \ref{th:3.3.2}, 
$X$ is countably paracompact.

Now to see that $X$ is $\lambda$-shrinkable, 
let $\mathscr{U} = \{U_\alpha: \alpha < \lambda\}$ be an open cover of $X$. By \textsc{Claim 1}, 
for each $m, n < \omega$, there exists a family 
$\mathscr{V}_{m,n} = \{V_{m,n,\alpha}: \alpha< \lambda\}$ of open sets of $X$ such that 
$\overline{V}_{m,n,\alpha} \subset U_\alpha$ for each $\alpha<\lambda$ and 
$\mathscr{V} = \bigcup_{m,n \in \omega} \mathscr{V}_{m,n}$ covers $X$. By \textsc{Claim 2},
 there exists a locally finite open cover $\{W_{m,n}: m, n < \omega\}$ such that 
 $W_{m,n} \subset \bigcup\mathscr{V}_{m,n}$. 
 For each $\alpha < \lambda$, let 
 $H_\alpha = \bigcup_{m,n < \omega} (W_{m,n} \cap V_{m,n,\alpha})$
It follows that $\mathscr{H} = \{H_\alpha\}$ is an open cover of $X$ such that 
$\overline{H}_\alpha \subset U_\alpha$ for each $\alpha$, so $X$ is shrinkable.
\end{proof}

{\corollary{} \label{cor:5.2.5}
 $X$ is countably shrinkable iff $X$ is normal and countably paracompact.
}

\begin{proof}
\((\Rightarrow)\) By Lemma \ref{lem:5.2.2}. \((\Leftarrow)\) By Theorem \ref{th:5.2.4}.
\end{proof}

{\lemma{} \label{lem:5.2.6}
If every open cover $\mathscr{U} = \{U_\alpha: \alpha\in \Delta\}$ of $X$ has an open refinement
$\mathscr{V} = \{V_{n, \alpha}: n < \omega, \alpha\in\Delta\}$ such that 
$\overline{V}_\alpha \subset U_\alpha$ for each $n, \alpha$, then
$X$ is normal.
}

\begin{proof}
Let \(A\) and \(B\) be disjoint closed sets in \(X\). Then the open cover 
\(\{ X \setminus A, X \setminus B \}\) has an open refinement 
\(\mathscr{V} = \{ V_{n,0}, V_{n,1} : n < \omega \}\) such that 
\(\overline{V}_{n,0} \subset X \setminus B\) and
 \(\overline{V}_{n,1} \subset X \setminus A\) for each \(n < \omega\).
For each \(k < \omega\), let 
$$W_{2k} = V_{k,0}, \quad W_{2k+1} = V_{k,1}.$$
Then \(\mathscr{W} = \{ W_n : n < \omega \}\) is an open cover of \(X\). 
For \(n < \omega\), let 
\(B_0 = W_0, B_n = W_n \setminus \overline{\bigcup_{i<n}W_i}\).
Then \(\mathscr{B} = \{ B_n : n < \omega \}\) is a locally finite refinement of \(\mathscr{W}\).
Let $F = \overline{\bigcup_{k<\omega} B_{2k}}, \quad K = \overline{\bigcup_{k<\omega} B_{2k+1}}.$
Then  \(F\) and \(K\) are closed sets such that 
$$A \subset X \setminus F, \quad
B \subset X \setminus K, \quad
(X \setminus F) \cap (X \setminus K) = \emptyset.$$
Thus \(X\) is normal.
\end{proof}


{\theorem{} \label{the:5.2.7}
For any space \(X\), the following are equivalent:

\textup{(i)} \(X\) is \(\lambda\)-shrinkable.

\textup{(ii)} Every open \(\lambda\)-cover \(\mathscr{U}=\{U_\alpha: \alpha < \lambda\}\) of \(X\) 
has a functionally open refinement \(\mathscr{V} = \{ V_{n,\alpha} : n < \omega, \alpha < \lambda \}\) 
such that \(\overline{V}_{n,\alpha} \subset U_\alpha\) for $n < \omega, \alpha < \lambda$.

\textup{(iii)} Every open \(\lambda\)-cover \(\mathscr{U}=\{U_\alpha: \alpha < \lambda\}\) of \(X\) 
has an open refinement \(\mathscr{V} = \{ V_{n,\alpha} : n < \omega, \alpha < \lambda \}\) 
such that \(\overline{V}_{n,\alpha} \subset U_\alpha\) for $n < \omega, \alpha < \lambda$.

\textup{(iv)} \(X\) is normal and every open \(\lambda\)-cover  
\(\mathscr{U}=\{U_\alpha: \alpha < \lambda\}\) of \(X\) 
has a closed refinement \(\mathscr{B} = \{ B_{n,\alpha} : n < \omega, \alpha < \lambda \}\) 
such that \(B_{n,\alpha} \subset U_\alpha\) for $n < \omega, \alpha < \lambda$.
}

\begin{proof}
(i) \(\Rightarrow\) (ii). By Lemma \ref{lem:5.2.2}.
			
(ii) \(\Rightarrow\) (iii). Clear.
			
(iii) \(\Rightarrow\) (iv). By Lemma \ref{lem:5.2.6}.
			
(iv) \(\Rightarrow\) (i). Assume (iv). By Theorem \ref{th:5.2.4}, it suffices to prove that \(X\) is 
\(\lambda\)-$C\omega^2$-refinable. 
Let \(\mathscr{U} = \{ U_\alpha : \alpha < \lambda \}\) be an open cover of \(X\). 
Let \(\{ B_{n,\alpha} : n < \omega, \alpha < \lambda \}\) be a closed cover of \(X\) 
such that \(B_{n,\alpha} \subset U_\alpha\) for each \(n < \omega, \alpha < \lambda\).
For \(i < \omega\), let \(\mathscr{B}_i = \{ B_{i,\alpha} : \alpha < \lambda \}\).
			
\textsc{Claim.} For each \(m < \omega\), there exists a family \(\mathscr{F}_{m,0} = \{ F_{m,0,\alpha} : \alpha < \lambda \}\) in $X$ such that
\begin{enumerate}
\item[(1)] \(F_{m,0,\alpha} \subset U_\alpha\) for each \(\alpha < \lambda\),
\item[(2)] \(\mathscr{F}_{m,0}\) is a family of closed subsets 
 of \(X\setminus\bigcup_{i<m}(\bigcup\mathscr{F}_{i,0})\), and
\item[(3)] \(\bigcup_{i \le m} (\bigcup \mathscr{F}_{i,0}) = \bigcup_{i \le m} (\bigcup \mathscr{B}_i)\).
\end{enumerate}

\textsc{Proof of Claim}.
We use induction on $m$.
Let $\mathscr{F}_{0,0} = \{\lambda F_{0,0, \alpha}: \alpha < \lambda$ 
where $F_{0, 0, \alpha} = B_{0,\alpha}$. Let $m < \omega$ and assume that the families 
$\mathscr{F}_{i,0} = \{F_{i,0,\alpha}: \alpha < \lambda\}$ have been 
constructed such that the conditions are satisfied for all $i \le m$. For $m+1$ let
$F_{m+1,0,\alpha} = B_{m,n,\alpha} - \bigcup_{i \le m} (\bigcup \mathscr{F}_{i,0})$.
Then $\mathscr{F}_{m+1,0} = \{F_{m+1,0,\alpha}\}$ is a family of closed subsets of $X$, and $F_{m+1,0,\alpha} \subset V_\alpha$ for each $\alpha$. Then
$\bigcup \mathscr{F}_{m+1,0} = \bigcup \mathscr{B}_{m,n} - \bigcup_{i \le m} (\bigcup \mathscr{F}_{i,0}) = \bigcup \mathscr{B}_{m,n} - \bigcup_{i \le m} (\bigcup \mathscr{F}_{i,0})$.
Thus
$\bigcup_{i \le m+1} (\bigcup \mathscr{F}_{i,0}) = \bigcup_{i \le m} (\bigcup \mathscr{F}_{i,0}) \cup (\bigcup \mathscr{F}_{m+1,0}) = \bigcup_{i \le m+1} (\bigcup \mathscr{B}_{i,0})$.
It follows that the family $\mathscr{F}_{m,n,0}$ satisfies the condition (1)–(3) above. This completes the proof of the claim.
By (3), $\bigcup (\bigcup \mathscr{F}_{m,0}) = X$. For each $m < \omega$ and $n$, let $\mathscr{S}_{m,n} = \{\phi\}$. It follows from claim that $\mathscr{S} = \bigcup_{m,n < \omega} \mathscr{S}_{m,n}$ is a $\sigma$-closure refinement of $\mathscr{U}$, so $X$ is $\sigma$-closure-refinable.
\end{proof}


{\theorem{} \label{the:5.2.8}
For any $X$, the following are equivalent:
(i) $X$ is normal.
(ii) Every $\sigma$-cover of $X$ has a functionally open shrink.
(iii) Every $\sigma$-cover of $X$ has an open shrink.
(iv) Every $\sigma$-cover of $X$ has a finite shrinkable.
(v) Every $\sigma$-cover of $X$ has a closed shrink.
}

\begin{proof}
 (i) $\implies$ (ii). Suppose $\mathscr{U}$ is a $\sigma$-cover of $X$. 
 Let $\mathscr{U} = \bigcup_{m,n < \omega} \mathscr{U}_{m,n}$ 
 be a refinement of $\mathscr{U}$.
\end{proof}

Here is the text extracted from the document:

where each $V_{m,n,\alpha} = F_{m,n,\alpha}$ with $U_{m,n}$ and $U_{m,n}$ are $F_{m,n,\alpha}$ and each $m,n,\alpha$ closed subsets of $X$.
Claim. For each $m,n$ there is a functionally open family $\mathcal{V}_{m,n} = \{V_{m,n,\alpha}\}$ such that:
(1) $\overline{V}_{m,n,\alpha} \subset U_\alpha$ for each $\alpha$,
(2) $\mathscr{V}_{m,n}$ is functionally open, and
(3) $\bigcup_{(p,q)<(m,n)} (\bigcup \mathscr{V}_{p,q}) < \bigcup_{(p,q)\le(m,n)} (\bigcup \mathscr{V}_{p,q})$.
Proof. We use induction on $(m,n)$. Suppose $\mathscr{S}_{0,0} = \{\emptyset\}$ and let $\dots$
Let $m,n \in \omega$ and assume that the families $\mathscr{V}_{p,q}$ have been constructed such that the conditions (1)–(3) are satisfied for all $(p,q) < (m,n)$. Then
$\langle p,q \rangle < (m,n) (\bigcup \mathscr{V}_{p,q}) < \bigcup_{(p,q)\le(m,n)} (\bigcup \mathscr{V}_{p,q})$
$L = \bigcup_{(p,q)<(m,n)} (\bigcup \mathscr{V}_{p,q})$
Then $\mathscr{F}_{m,n}(x-1)$ is a closed family in $X$. For each $\alpha$ there is a functionally open set $H_{m,n,\alpha}$ such that
$F_{m,n,\alpha} - L < H_{m,n,\alpha} \subset \overline{H}_{m,n,\alpha} < U_\alpha$.
Then
$\bigcup \mathscr{F}_{m,n} - L \subset \bigcup_{\alpha \in \mathscr{A}} H_{m,n,\alpha}$. Since $X - L$ is a closed subset of $X$, there is a functionally open set $K$ such that
$\bigcup \mathscr{F}_{m,n} - L \subset K \subset \overline{K} \subset \bigcup_{\alpha \in \mathscr{A}} H_{m,n,\alpha}$.
Let $V_{m,n,\alpha} = \dots$. Then $\mathscr{V}_{m,n} = \{V_{m,n,\alpha}\}$ is a functionally open family with $\overline{V}_{m,n,\alpha} \subset U_\alpha$ and such that
$\bigcup \mathscr{V}_{m,n} = K$ is functionally open. Since the set
$\bigcup \mathscr{F}_{m,n} - L \subset \bigcup \mathscr{V}_{m,n}$,

Thus $U_{m,n}$ is the desired functionally open family and the claim is proved.

By the claim and Theorem 2016, there is a LF functionally open cover $\{W_{m,n}\}$ for each $m,n$. Such that $W_{m,n} \subset \bigcup \mathscr{V}_{m,n}$.

For $\alpha$, let $G_\alpha = \bigcup (W_{m,n} \cap V_{m,n,\alpha})$. Then $\overline{G}_\alpha \subset U_\alpha$ for each $\alpha$, so $\mathscr{G}$ is an open shrink of $\mathscr{U}$.

(ii) $\implies$ (iii) $\implies$ (iv) $\implies$ (v): Clear.

(v) $\implies$ (i): Routine.

15.2.9. Corollary. For any $X$, the following are equivalent:

$X$ is normal.


(ii) (Price and smooth...). Every open cover has a closed shrink.


(iii) Every $\sigma$-disjoint refinement also has a closed refinement.


(iv) (Lefschetz 1942). Every $\sigma$-finite open cover of $X$ has an open shrink.



Proof. Since the ordinals are isomorphic well-ordered sets, any $\sigma$-disjoint refinement is a $\sigma$-closure refinement.

5.2.10. Lemma. Any normal countably metacompact space is countably paracompact.

73
{\corollary{\label{cor:5.2.9}}
For
}
17

Proof. Suppose $X$ is normal countably metacompact and $\mathscr{U} = \{U_n\}$ be an increasing open cover of $X$. Then $\mathscr{U}$ has a point-finite open refinement $\mathscr{V}$. By corollary 6.2.10, there is an open cover $\mathscr{W}$ of such that $\overline{W}_n \subset U_n$. By Theorem 4.9, $X$ is countably paracompact.

5.2.11. Theorem. For any $X$, the following are equivalent:

(i) $X$ is normal.


(ii) Every LF open cover of $X$ is normal.


(iii) $X$ is fine fully normal.



Proof. (i) $\implies$ (ii). Suppose $X$ is normal and $\mathscr{U}$ is a LF open cover of $X$. By corollary 5.2.9, $\mathscr{U}$ has a functionally open shrink $\mathscr{V}$. Then $\mathscr{V}$ is also LF. By Theorem 2.1.6, $X$ is normal.

(ii) $\implies$ (iii). Clear.

(iii) $\implies$ (i). Suppose $A$ and $B$ are disjoint closed sets in $X$. By (iii), the open cover $\{X-A, X-B\}$ has a point-star open refinement $\mathscr{S}$. Then $A \subset \bigcup \{S \in \mathscr{S} : S \cap A \neq \emptyset\}$ and $B \subset \bigcup \{S \in \mathscr{S} : S \cap B \neq \emptyset\}$, so $X$ is normal.

74
\end{document}
